%% kapa: exam=2 round=30 year=2019 period=3 subject=compensation_law subject_ko=감정평가 및 보상법규 seq=4 number=4 points=10
%% source: Q-Net 감정평가사 기출문제 (cst00309) artlSeq=5205595 file=raw/1008/5205595/감정평가 및 보상법규.pdf
%% conversion: kordoc 4.24.51 -> tools/gen2.py (auto); review=proofread-gpt-terra; flags=none
\begin{longquestion}{【문제 4】}
「공익사업을 위한 토지 등의 취득 및 보상에 관한 법률」제26조는 수용재결 신청 전에 사업시행자로 하여금 수용대상 토지에 관하여 권리를 취득하거나 소멸시키기 위하여 토지소유자 및 관계인과 교섭하도록 하는 협의제도를 규정하고 있다. 이에 따른 협의가 수용재결 신청 전의 필요적 전치절차인지 여부와 관할 토지수용위원회에 의한 협의성립의 확인의 법적 효과를 설명하시오. (10점) \par
\end{longquestion}
